\section{Evaluation}
\subsection{Configurations and Measurement Protocol}
We compare six handoff configurations in serial and double-buffered execution: faithful full image (allocated-region clearing and transfer), lean full image (used-region clearing and transfer), CPU-generated delta, CPU-generated delta + on-chip base reuse, hardware-generated delta (ReasonLink), and hardware-generated delta + on-chip base reuse. A thirteenth configuration replays classifier logits with the faithful serial path. The replay is a diagnostic that removes actual NVDLA execution; it is not a measured faster neural accelerator.

The reported dataset contains one completed 100-puzzle run per configuration. Failed attempts are excluded from these completed-run timings and discussed separately. Model, encoding, predicted digits, and solver settings are held fixed. We report whole-run wall time divided by puzzle count and medians of instrumented per-puzzle stages. The median of a stage and the mean wall time are different statistics and should not be added or subtracted as if they were a single representative execution.

The current wall-time interval excludes initial model, dataset, and base-image loading. Each run has zero warm-up puzzles; the first solve loads on-chip base state even when subsequent solves reuse it. In the hardware-generated path, the inference interval also includes descriptor submission and credit waits. CPU-generated updates are computed after inference for the puzzle, whereas hardware-generated updates can overlap subsequent cell classifications.

\draftnote{Fix the workloads, correctness checks, randomized configuration order, and experiment budget. Repeat clean runs without SSH monitoring; retain raw and summarized data and report uncertainty across independent runs. Specify the warm-up policy and report cold-start costs separately from steady-state performance. Match CPU and hardware update payloads and generation schedules to isolate computation placement, and use consistent timing boundaries for conversion, submission, transfer, and synchronization.}

Some runs included SSH monitoring that competed for the board's slow CPU. The existing summarizer estimates excess inference time using each run's median inference time per given. For puzzle $i$, with $g_i$ givens, it computes
\begin{align}
b &= \operatorname{median}_i(T_{\mathrm{nn},i}/g_i),\\
\widehat{S} &= \sum_i\max(T_{\mathrm{nn},i}-b g_i,0),\\
\widehat{T}_{\mathrm{wall}} &= T_{\mathrm{wall}}-\widehat{S}.
\end{align}
This adjustment attributes all above-median inference excess to interference and may remove real variation as well. We show it as an estimate alongside measured wall time. It does not replace clean repeats or support statistical significance. A later five-configuration double-buffered batch ran without SSH monitoring and completed 500 puzzles without a hang; independent repeated batches are still needed.

\draftnote{Use clean repeated measurements for final headline comparisons. Separate NVDLA device time from runtime/driver overhead, and measure sensitivity to neural production rate, solver difficulty, and symbolic-state size. Compare against the optimized CPU-generated delta path as well as full-image baselines. Add a figure showing the measured regimes where CPU-generated deltas, hardware-generated deltas, and on-chip base reuse are beneficial.}

\begin{figure*}[t]
\centering
\begin{tikzpicture}
\begin{axis}[name=serial,width=0.49\textwidth,height=5.2cm,
 ybar,bar width=7pt,ymin=0,ymax=4.6,
 ylabel={Seconds per puzzle},title={Serial},
 xtick=data,xticklabels={Faithful,Lean,CPU delta,CPU delta+B,HW delta,HW delta+B},
 x tick label style={font=\scriptsize,rotate=25,anchor=east},
 enlarge x limits=0.12,grid=major,grid style={gray!15},
 legend style={font=\scriptsize,at={(0.98,0.98)},anchor=north east,legend columns=1},
 tick label style={font=\scriptsize}]
\addplot[fill=blue!55,draw=blue!70] table[x=index,y=measured_s,col sep=tab]{figures/serial-walltime.tsv};
\addplot[fill=orange!65,draw=orange!80!black] table[x=index,y=adjusted_s,col sep=tab]{figures/serial-walltime.tsv};
\legend{Measured,Adjusted estimate}
\end{axis}
\begin{axis}[at={(serial.east)},anchor=west,xshift=0.6cm,
 width=0.49\textwidth,height=5.2cm,ybar,bar width=7pt,ymin=0,ymax=4.6,
 title={Double-buffered},
 xtick=data,xticklabels={Faithful,Lean,CPU delta,CPU delta+B,HW delta,HW delta+B},
 x tick label style={font=\scriptsize,rotate=25,anchor=east},
 enlarge x limits=0.12,grid=major,grid style={gray!15},
 tick label style={font=\scriptsize}]
\addplot[fill=blue!55,draw=blue!70] table[x=index,y=measured_s,col sep=tab]{figures/double-walltime.tsv};
\addplot[fill=orange!65,draw=orange!80!black] table[x=index,y=adjusted_s,col sep=tab]{figures/double-walltime.tsv};
\end{axis}
\end{tikzpicture}

\caption{Mean wall seconds per puzzle from the current 100-puzzle board runs. Measured wall times and the repository's interference-adjusted estimates are shown separately. Each bar represents one completed run; no error bars or significance claims are available. CPU delta and HW delta denote CPU-generated delta and hardware-generated delta (ReasonLink); +B enables on-chip base reuse for either path.}
\label{fig:wall}
\end{figure*}

\subsection{H1: Full-Image Preparation Is Substantial}
Table~\ref{tab:serial} reports serial-stage medians. The faithful baseline spends approximately 722\,ms building its image and 609\,ms in transfer, compared with approximately 11\,ms in SAT-Accel loading and search. Its median instrumented handoff fraction is 34.9\%. Lean clearing and transfer reduce these stages to approximately 412 and 73\,ms, respectively, but preparation remains visible at a 16.3\% handoff fraction.

CPU-generated deltas reduce conversion and transfer to approximately 0.148 and 0.336\,ms. The CPU-generated delta path's median handoff fraction is 0.042\%; the hardware-generated delta path's recorded fraction is approximately 0.024\%. ReasonLink descriptor pushes and credit waits occur inside the inference interval, so this fraction is not a complete hardware-interface cost. Separate diagnostic timing records approximately 18\,$\mu$s of CPU push and credit work per given, or about 0.6\,ms per puzzle.

Host-side transfer volume falls from 5,293,556 bytes in the faithful path and 632,436 bytes in the lean path to 3,940 bytes for CPU-generated delta and 2,948 bytes for hardware-generated delta. These counts exclude NVDLA traffic and SAT-Accel's internal copy-in reads. They therefore quantify the handoff update protocol, not total system memory traffic.

\begin{table*}[t]
\centering
\caption{Serial board-stage medians in milliseconds per puzzle. Base reuse denotes on-chip base reuse; hardware-generated delta is ReasonLink. Host update bytes are distinct from internal solver loads. ReasonLink submission and credit handling are included in $T_{\mathrm{nn}}$; zero CPU conversion does not imply zero interface work.}
\label{tab:serial}
\small
\begin{tabular}{@{}lcrrrrrr@{}}
\toprule
Delta generation / baseline & Base reuse & $T_{\mathrm{nn}}$ & $T_{\mathrm{conv}}$ & $T_{\mathrm{xfer}}$ & $T_{\mathrm{sync}}$ & $T_{\mathrm{sat}}$ & Update bytes\\
\midrule
Faithful full image & No & 2537.826 & 721.970 & 609.000 & 0.723 & 10.995 & 5,293,556 \\
Lean full image & No & 2530.492 & 411.617 & 73.198 & 0.551 & 10.645 & 632,436 \\
CPU-generated delta & No & 2547.761 & 0.148 & 0.336 & 0.563 & 10.192 & 3,940 \\
CPU-generated delta & Yes & 2458.728 & 0.147 & 0.336 & 0.545 & 3.058 & 3,940 \\
Hardware-generated delta & No & 2618.511 & 0.000 & 0.060 & 0.552 & 10.118 & 2,948 \\
Hardware-generated delta & Yes & 2615.484 & 0.000 & 0.060 & 0.538 & 3.005 & 2,948 \\
\bottomrule

\end{tabular}
\end{table*}

\subsection{CPU-Generated Versus Hardware-Generated Deltas}
Figure~\ref{fig:wall} shows that the optimized paths cluster near 2.5\,s per puzzle after adjustment. CPU-generated delta takes \SerialDeltaAdjusted\,s per puzzle and hardware-generated delta takes \SerialRLAdjusted\,s, a difference of only \DeltaRLDifference\% in the adjusted estimates. With on-chip base reuse, the corresponding values are \SerialDeltaResidentAdjusted\,s and \SerialRLResidentAdjusted\,s. Given the unequal monitoring interference and absence of independent repeats, these differences do not establish a hardware speedup.

Relative to faithful serial execution, ReasonLink's adjusted comparison is approximately \FaithfulRLSpeedup$\times$; relative to lean serial execution it is approximately \LeanRLSpeedup$\times$. These gains combine the representation change and execution mechanisms. The CPU-generated delta result shows that base reuse and compact updates already capture almost all of the available end-to-end benefit in the current platform.

\subsection{On-Chip Base Reuse Reduces Solver Loading}
Enabling residency in the CPU-generated delta path reduces median serial SAT time from \DeltaSat\,ms to \DeltaResidentSat\,ms, approximately \ResidentReduction\%. The hardware-generated delta path shows a similar change, from approximately 10.1 to 3.0\,ms. This stage improvement has little effect on the approximately 2.5-second pipeline because inference dominates.

RTL sequences provide separate evidence for the mechanism. After a common fresh first solve, resident puzzle 3 takes 100,553 busy cycles versus 207,383 with full copy-in; resident puzzle 5 takes 81,549 versus 188,379. Read beats fall from approximately 81,400 to 20,800. These are testbench results with a 30-cycle read latency and one outstanding burst, rather than predictions of board runtime. They establish reduced loading under that memory model, with unchanged answers and search counts.

\subsection{H2: Overlap Does Not Remove Repeated Preparation}
Faithful double buffering measures approximately 3.944\,s per puzzle, versus 3.978\,s serial; lean double buffering measures 2.962\,s versus 3.037\,s serial. The current overlap arrangement therefore leaves substantial preparation exposed for full-image paths. It does not reduce the host update bytes.

CPU-generated and hardware-generated delta double-buffered runs measure approximately 2.47\,s per puzzle. Differences from serial include overlap, unequal interference, and a roughly 1\% change in observed per-cell neural time in the later batch. The present data cannot attribute all throughput differences to double buffering. The optimized paths yield slightly negative $R$, consistent with overlap between solver and neural execution.

\subsection{Correctness and Boundary Cases}
All 13 accepted configurations agree on SAT/UNSAT and grids for all 100 board puzzles. The summarizer reports 98 reference-equivalent cases and two valid alternative grids per configuration. Its reference-equivalent category includes matching UNSAT outcomes; the 98 count is not a count of 98 identical solved grids. Alternative grids satisfy Sudoku constraints and retain every predicted given. These checks establish interface equivalence, not 100\% recovery of ground-truth puzzles. Across the full 1,000-puzzle emulated workload, the recorded ground-truth puzzle accuracy is 79.0\%.

Functional resident-state validation covers 1,000 Sudoku puzzles and 192 fixed-decision benchmark instances. Resident and fresh executions agree on models and search counters, and the restored base regions match byte for byte. Without tail cleanup, stale literal-store bytes remain in some instances, demonstrating why explicit restoration is needed.

NVDLA and the emulator differ in the logits of 925--926 given-cell classifications per board run, but the selected digits agree on the evaluated 100 puzzles. Logit-level bit equivalence has therefore not been established for the retrained board model. Two earlier NVDLA mid-task hangs also required FPGA reprogramming; their trigger is unresolved. The later uninterrupted batch is evidence of a successful run, not proof that the reliability issue has been eliminated.

\draftnote{Quantify logit differences, measure ground-truth board accuracy separately, and report failed-attempt frequency and timeout policy. Add a controlled failure/boundary study beyond these observed integration cases.}
